\section{ADMM --- sparsity by soft thresholding}
\label{sec:admm}

\paragraph{How it works.}
ADMM is the MAP estimator for a \emph{sparse}, non-negative object under Gaussian noise: its
prior is hard-wired ($\ell_1$ + positivity), so it ignores $R$, $\lambda$, $\kappa$. It alternates
a data step, a soft threshold and a dual update:
\begin{quote}\small
$t = \kappa\, \max(g)/s_1$ (threshold);\quad
$u \leftarrow (H^{\!\top}\!H+\mu I)^{-1}(H^{\!\top} g + \mu(f-\eta))$;\quad
$f \leftarrow \max(u+\eta-t,\,0)$;\quad $\eta \leftarrow \eta + \mu(u-f)$.
\end{quote}
The Laplacian prior $\exp(-\tau\lVert f\rVert_1)$ says ``most voxels are ${\approx}0$, a few are
bright'' --- the right model for point-like / filamentary structures, the wrong one for a
continuous membrane.

\paragraph{Parameters.}
Two knobs move the solution, and awkwardly. The threshold ratio $\kappa$ (\texttt{threshold\_ratio})
sets the sparsity, but its useful band cannot be predicted a priori --- the data step spreads the
object across the ${\sim}47$ null-space directions, so every voxel of the ridge estimate is small
and a diffuse truth can be emptied by a modest $\kappa$. The penalty $\mu$ carries \emph{three}
coupled roles at once (data-step ridge, prior weight $\tau=\mu\kappa\max(g)/s_1$, and dual step),
and the dual step caps it: $\mu > 1.618$ oscillates or diverges (measured: $\mu=2.5\to$ NaN).
\texttt{iter}$=200$ is a generous budget (the default 20 is under-iterated).

\paragraph{Standard use.}
$\kappa=0.1$, $\mu=0.5$ ($\approx s_3^2$, inside the stable band). On a concentrated object at
low noise this is already close to the best ADMM can do.

\paragraph{Optimal use.}
Sweep $\kappa$ log-spaced (here $\kappa\in\{0.1,\,0.01\}$) and raise it with the noise; keep
$\mu \in [s_3^2,\,1.6] = [0.3,\,1.6]$. The measured standard-vs-optimal figures:
\IfFileExists{tables/algo_admm.tex}{\paragraph{matirf}
\begin{center}\begin{adjustbox}{max width=\linewidth}
\begin{tabular}{lllllllll}
\toprule
use & dataset & noise & params & NMSE & PSNR & Depth Error (nm) & Stack Recovery & chi2\_ratio \\
\midrule
standard / optimal & cell & 0 & kappa=0.1, mu=0.5 & 0.0318 & 28.8 & 4 & 0.636 & 5.2e-08 \\
standard / optimal & cell & 0.02 & kappa=0.1, mu=0.5 & 0.619 & 15.9 & 41 & 0.152 & 0.604 \\
standard / optimal & cell & 0.05 & kappa=0.1, mu=0.5 & 0.899 & 14.3 & 68 & 0.0758 & 0.589 \\
optimal & esoubies & native & kappa=0.1, mu=0.5 & — & — & — & — & 19.2 \\
standard & fibres & 0 & kappa=0.1, mu=0.5 & 0.251 & 30.1 & 15 & 0.336 & 0.000278 \\
optimal & fibres & 0 & kappa=0.01, mu=0.5 & 0.0406 & 38.0 & 10 & 0.966 & 4.86e-06 \\
standard / optimal & fibres & 0.02 & kappa=0.1, mu=0.5 & 0.882 & 24.6 & 64 & 0.0517 & 0.436 \\
standard / optimal & fibres & 0.05 & kappa=0.1, mu=0.5 & 0.973 & 24.2 & 95 & 0.0172 & 0.411 \\
standard & vesicles & 0 & kappa=0.1, mu=0.5 & 0.142 & 34.6 & 8 & 0 & 0.00024 \\
optimal & vesicles & 0 & kappa=0.01, mu=0.5 & 0.021 & 42.9 & 8 & 1 & 5.3e-06 \\
standard / optimal & vesicles & 0.02 & kappa=0.1, mu=0.5 & 0.848 & 26.8 & 66 & 0 & 0.413 \\
standard / optimal & vesicles & 0.05 & kappa=0.1, mu=0.5 & 0.974 & 26.2 & 94 & 0 & 0.39 \\
\bottomrule
\end{tabular}\end{adjustbox}\end{center}
}%
  {\emph{[pending the run]}}

\IfFileExists{figures/matirf_vesicles_ADMM_noise00.png}{%
\begin{figure}[H]\centering
  \includegraphics[width=.7\linewidth]{matirf_vesicles_ADMM_noise00}
  \caption{ADMM on \texttt{vesicles} at zero noise --- where the sparsity prior localises the
  point-like structure well.}
\end{figure}}{}

\paragraph{When it works, when it fails.}
ADMM's wall is its prior. On a diffuse or continuous object $\ell_1$ is the wrong model, so there
is little-to-no useful $\kappa$ band: $\kappa$ too large empties the image (the dominant failure,
and it arrives early), $\kappa\approx0$ is noisy and lets \texttt{iter} become an accidental
regulariser. Concentrated truths (vesicles, filaments) have a band, but a narrow one. Use ADMM
for sparse objects at low noise; do not expect it to reconstruct a membrane.
